\subsection{Replacing the blind terminal slice in the exact rank-two pair}
\label{sec:exact-rank-recovery}

\begin{theorem}
\label{th:exact-rank-recovery}
Take the pair of Theorem~\ref{th:rank-two}
with $gq\ne0$ and a single-mode seed at nodes $a<b$.
Write its phase as $w_\epsilon=\epsilon u_{k_0}(t)\cos(k_0\theta)$.
There is an explicit rational delay $d=\pi/L$ with
$0<d<b-a$ such that, for all sufficiently small
nonzero $|\epsilon|$, the full orbit fields at
$c=a+d$ give a locally bi-Lipschitz missing-$P_t(a)$
inverse around either rank-two history.
For each integer $s\ge1$, the snapshot is measured
in $H^{s+1}$ and the velocity in $H^s$.
The lower bound for the derivative can be taken
to be a positive constant times $|\epsilon|$.
\end{theorem}

\begin{proof}
Choose an integer $L>2k_0$ satisfying the short-gap
conditions of Proposition~\ref{prop:safe-gap}, and
large enough that $u_{k_0}(c)\ne0$.
Let $\mathcal P_\epsilon$ and $\mathcal Q_\epsilon$
denote the two linearized terminal observation
operators in the scaled orthonormal frame of
\eqref{eq:transport-leading}. Their dependence on
$\epsilon$ in the operator norm $H^s\to H^{s+1}$
is analytic on a small interval.
At $\epsilon=0$ the transformed background is a
vertical homogeneous solution with $Q$ constant.
The $P$-variation then solves the scalar
Euler--Poisson--Darboux equation and the $Q$-variation
vanishes for the given missing-velocity direction.
Moreover, the first derivative of the $P$-observation
in $\epsilon$ vanishes. In the coordinate linearized
$P$ equation the terms beyond the free damped wave
are proportional to $Q_t^2-Q_\theta^2$ times
$\delta P$, or to a first derivative of $Q$ times
a first derivative of $\delta Q$.
Both are quadratic in the amplitude at this
reference. The time rescaling of $P$ is independent
of $\epsilon$. Hence
\[
 \mathcal P_\epsilon=\mathcal P_0+O(\epsilon^2),
 \qquad
 \mathcal Q_\epsilon=\epsilon\mathcal Q_1+O(\epsilon^2)
\]
in the stated operator norms.

On the multiples of $L$, the two shifts of the
zero-mean primitive $h=\partial_\theta^{-1}f$ in
\eqref{eq:transport-leading} agree, with sign
$(-1)^{k/L}$ on mode $k$.
The rotation angle has first derivative
$q\,u_{k_0}(c)\cos(k_0\theta)$ at $\epsilon=0$,
because $B'(0)=q$ for this normalization.
Thus on the nonzero multiples of $L$ the leading
part of $\mathcal Q_1$ is
\[
 f\longmapsto q\,u_{k_0}(c)\cos(k_0\theta)\,
                 \mathcal J\,\partial_\theta^{-1}f ,
\]
where $\mathcal J$ is the sign operator
$(-1)^{k/L}$.
The differentiated remainder is bounded
from $H^{s-1}$ to $H^{s+1}$, as follows by
differentiating the smooth transport coefficients
and the forced remainder equation.
Multiplication by $\cos(k_0\theta)$ sends each
multiple of $L$ to two disjoint sidebands
$k\pm k_0$. They do not overlap with sidebands
of any other multiple of $L$ because $L>2k_0$.
The two shifted Sobolev weights are comparable to
the unshifted ones. The displayed leading operator
therefore has a uniform positive lower bound.
After discarding finitely many resonant modes,
the compact remainder has arbitrarily small norm
on the remaining high resonant subspace.
Consequently $\mathcal Q_1$ is bounded below there.

Split $f=f_G+f_B$ into the complementary modes
(the nonmultiples and finitely many small
multiples of $L$) and the remaining large
multiples. The scalar short-gap transfer
$\mathcal P_0$ is bounded below on $G$ by the
nonresonant free phases and the nonzero finite
transfers of Proposition~\ref{prop:safe-gap}.
Projecting the $P$ record onto $G$ thus gives
\[
 \|f_G\|_{H^s}\le C\big(
    \|\mathcal P_\epsilon f\|_{H^{s+1}}+
      \epsilon^2\|f\|_{H^s}\big).
\]
The first-order estimate on $B$, together with
the triangle inequality for $\mathcal Q_1 f_G$,
gives
\[
 \|f_B\|_{H^s}\le C'\left(
   \frac{\|\mathcal Q_\epsilon f\|_{H^{s+1}}}{|\epsilon|}
       +\|f_G\|_{H^s}+|\epsilon|\|f\|_{H^s}\right).
\]
Combining and absorbing the small-amplitude terms
bounds the inverse by $O(|\epsilon|^{-1})$.
The constants may depend on the selected times
and the seed parameters.
The two sign branches use the same argument.
The local nonlinear inverse follows by continuity
of the Cauchy derivative as in
Theorem~\ref{th:fixed-rank-two}.
\end{proof}

The order of the lower gain is sharp as
$\epsilon\to0$ for the rational placement used here.
One can take arbitrarily high multiples of $L$.
The scalar transfer of $\mathcal P_0$ then tends
to zero in energy normalization, while the
uniform remainder estimates give contributions of
order $\epsilon^2$ and $|\epsilon|$ respectively
in the $P$ and $Q$ channels.

\subsection{Which extra terminal field supplies stable information}
\begin{proposition}
\label{prop:rank-two-fields}
For the exact small-amplitude rank-two pair in
Theorem~\ref{th:exact-rank-recovery}, the off-node
observation $P(c),\lambda(c)$ gives the same local
energy-norm stability even without a terminal $Q(c)$.
It is enough to read $\lambda_\theta(c)$ in $H^s$.
In contrast, no fixed finite collection of later $P$
snapshots alone gives an energy-normalized local
Lipschitz inverse at these backgrounds.
\end{proposition}

\begin{proof}
Write the two transport vectors in the frame as
$A_\pm=(\cos\rho,\mp\sin\rho)$.
The momentum constraint gives the leading
mean-zero part of the rescaled metric variation as
\[
 -\big(C_-H-C_+\mathcal T H\big),\qquad
 C_-=\langle A_-,U_t+U_\theta\rangle,\quad
 C_+=\langle A_+,U_t-U_\theta\rangle .
\]
Here the rescaling of $\delta\lambda$ is $\sqrt{c/a}/c$.
This follows by substituting the two leading orbit
transport terms in the momentum constraint and taking
a spatial primitive, as in Section~\ref{sec:metric-positive}.
All omitted terms map $H^{s-1}$ to $H^{s+1}$.
In the Ernst family the half difference of the
coefficients is
\[
 \frac{C_--C_+}{2}
 =\cos\rho\,\widetilde P_\theta
    +\sin\rho\,e^{\widetilde P}\widetilde Q_t
 =w_\theta[-\cos\rho\,P_w+\sin\rho\,C e^{-P(w)}]
 =-g w_\theta .
\]
The last equality holds because the seed geodesic
tangent with initial components $(g,q)$ is parallel:
its components at $w$ are obtained by a rotation
through $-B(w)=-\rho$.
At zero amplitude the linearized
$\lambda_\theta(c)$ equals $2\delta P_\theta(c)$.
Subtract this known channel and take the rescaled
mean-zero primitive of the difference.
Its linearization in amplitude, on modes that
are multiples of $L$, has the principal term
\[
 -2g k_0u_{k_0}(c)\sin(k_0\theta)\,
               \mathcal J\,\partial_\theta^{-1}f .
\]
This is bounded below on the high resonant
subspace by the same disjoint-sideband argument,
since $g\ne0$ and $L>2k_0$.
The remainder and the good-mode $P$ bounds are
the same as in Theorem~\ref{th:exact-rank-recovery}.
Thus $P$ and the metric channel give a bounded-below
derivative, and then a local nonlinear inverse.

For the negative statement, at each later time $t_i$
the leading $P$ component is a smooth multiplier
$\cos\rho_i$ times the difference of the two translated
primitives. The shifts are opposites and the $P$
components of their transported vectors are equal.
By simultaneous rational approximation choose a
sequence $k_m\to\infty$ with
$\max_i|\sin(k_m(t_i-a))|\to0$.
The corresponding unit-velocity Fourier modes in $H^s$
give leading $P$ norms tending to zero at every
chosen time. The compact remainders tend to zero
as well by \eqref{eq:transport-leading}.
A pairwise local Lipschitz inverse for these $P$
observations would imply a lower bound for the
derivative at the reference background, a
contradiction.
\end{proof}
